% Оформляем как главу без номера, добавляем в оглавление
\chapter*{ЗАКЛЮЧЕНИЕ}
\addcontentsline{toc}{chapter}{ЗАКЛЮЧЕНИЕ}

В ходе выполнения дипломного проекта была решена актуальная научно-техническая задача по разработке и реализации автоматизированной системы мониторинга параметров окружающей среды на базе технологий интернета вещей (IoT).

На основании проведенных исследований и проектных работ можно сделать следующие выводы:

\begin{enumerate}
    \item Проведен комплексный анализ существующих технических решений и протоколов передачи данных. Установлено, что для распределенных систем мониторинга наиболее эффективным является использование стека протоколов TCP/IP в сочетании с прикладным протоколом MQTT, что минимизирует объем передаваемого трафика.
    
    \item Обоснован выбор элементной базы: использование микроконтроллера ESP32 и цифрового датчика SHT31 позволило достичь высокой точности измерений температуры ($\pm 0,3$ $^\circ$C) и влажности ($\pm 2\%$) при сохранении низкой стоимости конечного устройства.
    
    \item Разработана принципиальная электрическая схема и топология печатной платы узла сбора данных. Применение современных методов фильтрации цепей питания и оптимизация трассировки обеспечили стабильную работу Wi-Fi модуля в условиях промышленных помех.
    
    \item Спроектированы и реализованы алгоритмы работы встроенного программного обеспечения, поддерживающие режим глубокого сна (Deep Sleep). Это позволило увеличить расчетный срок автономной работы устройства от одного заряда аккумулятора до 12 месяцев.
    
    \item В экономическом разделе подтверждена целесообразность внедрения разработки. Срок окупаемости проекта составил 1,78 года при значительном снижении рисков выхода из строя серверного оборудования.
    
    \item В разделе охраны труда предложены мероприятия по обеспечению электробезопасности и рассчитаны параметры освещенности рабочего места, соответствующие действующим нормам СН 2.04.03-2020.
\end{enumerate}

Разработанная система прошла опытную эксплуатацию, в ходе которой была подтверждена стабильность передачи данных и корректность работы алгоритмов обработки исключительных ситуаций. Результаты работы могут быть внедрены на предприятиях ИТ-сектора, а также в системах автоматизации жилых зданий.

Поставленные задачи выполнены в полном объеме, цель дипломного проекта достигнута.

\newpage